\clearpage\section{Document, retrieval and generation pipelines}
\subsection{Ingestion: read, chunk and embed}
The API checks the format and the streamed size, stores the file in object storage and creates a document and a job. The worker then reads, chunks and embeds the document as separate steps, each retried on its own, and deletes the file after embedding (Figure~\ref{fig:ingestion}).
\diagram{ingestion}{Asynchronous ingestion with bounded retries.}{ingestion}
\subsubsection{Readers and source locators}
PDF is read with PyMuPDF and layout rules, with OCR for pages without text and the formula and figure repairs of Section~\ref{sec:results-extraction}. PPTX keeps slide text, tables, SmartArt and speaker notes. Recordings are transcribed by Whisper in 600-second windows. Every reader returns units of text with a page, slide or timestamp locator.
\subsubsection{Chunking and embedding}\label{sec:chunking}
Chunking follows the structure of the source before applying size limits. Short pages or slides can share a chunk, but no chunk spans more than two of them, so a citation still names a page. A long page is split at word boundaries and keeps its locator with a part index. Transcript segments are accumulated with their start and end times. DOCX and HTML use recursive splitting at paragraphs, lines and sentences\citep{recursivesplit}. The file name and the section or slide heading are prefixed to each chunk. Table~\ref{tab:chunk-parameters} gives the settings; lengths are estimated tokens, not the embedding model's tokenizer.
\begin{table}[H]\centering\small\caption{Structure-aware chunking and generation-context parameters.}\label{tab:chunk-parameters}
\begin{tabularx}{\linewidth}{|p{4.3cm}|p{2.1cm}|Y|}\hline
Parameter & Value & Meaning\\\hline
Content budget & 340 & Estimated tokens before title/heading prefix\\\hline
Narrative overlap & 50 & Estimated-token overlap in recursive splitting\\\hline
Page/slide grouping & Target 250; at most 2 & Flush at target; at most two distinct pages/slides; check 340 before merging\\\hline
Long-page overlap & 40 & Word-boundary split of oversized page/slide units\\\hline
Transcript target & 150 & Timestamp accumulation target; long captions split with zero overlap\\\hline
Context prefix & Title + heading & Omit redundant heading; recompute estimated length\\\hline
Generation neighbors & $\pm 2$ & Same-document physical neighbors retain individual source codes\\\hline
\end{tabularx}\end{table}
Semantic chunking was not adopted: Qu et al. found no consistent gain that justifies its cost\citep{semanticchunking}. Embeddings come from the local Qwen3 model through the gateway alias \code{embed-default}; queries carry the retrieval instruction and documents do not. Changing the model export requires re-embedding every chunk, which \code{scripts/reembed_chunks.py} does.
\subsection{Retrieval}
\subsubsection{Scope, candidates and ranking}\label{sec:retrieval-policy}
Module retrieval is limited to the selected document, chat to the course. Dense cosine search and BM25 each return ranked candidates, and RRF with constant 60 fuses them\citep{bm25,rrf}. The BM25 index is cached per scope and rebuilt when the chunk set changes.

A reranker reorders the candidates: Jina multilingual INT8 by default, or the hosted Jev model when configured (Section~\ref{sec:rerank-choice}). LO and Content rerank 30 candidates and use adaptive-$k$ with floor eight. Practice/Quiz retrieval for each sub-objective reranks 15 candidates and uses fixed $k=4$. Module-edit chat and course Q\&A use fixed $k=8$. Adaptive selection finds the largest adjacent gap within the leading 90\%, adds five and enforces floor eight\citep{adaptive}. A flat score distribution retains the search window. Each process sends at most two concurrent requests to a self-hosted reranker and retries a busy one (429, 502--504, connection errors) up to five times with backoff, so a busy reranker no longer restarts the job. A failed Jev call keeps the hybrid order and is logged for alerting; in chat, any rerank failure after two attempts keeps the hybrid order.
\subsubsection{Physical evidence and corrective retrieval}
Up to two neighbouring chunks on each side of a hit are added for Content; each keeps its own id and source code. Before Content is written, a relevance check judges whether the passages cover the LO and sub-objectives; if not, it rewrites the query once, merges the new passages and judges again\citep{correctiverag}. The confidence and retry count are stored with the module.
\subsection{Generation orchestration}\label{sec:generation}
Figure~\ref{fig:generation} shows retrieval and generation flow with physical inline citations.
\diagram{retrieval-generation}{Retrieval and generation flow with physical inline citations.}{generation}
\subsubsection{Planner and learning objectives}
The planner extracts up to eight ranked concepts from the first 8,000 characters of the document and is cached per document and output language. A generation takes the top concept, retrieves passages for it and writes the terminal LO, its Bloom level and the sub-objectives. Concepts that appear only after the first 8,000 characters are not considered, a known limit.
\subsubsection{Content and inline citations}
Content has one heading per sub-objective. The prompt gives each passage a code, and important statements carry a marker after them, for example \code{Heap insertion takes O(log n). [S2]}. The server rejects unknown codes, derives the character span of each claim and resolves each code to its chunk and locator, which the source panel shows beside the lesson.
\needspace{.55\textheight}
\subsubsection{Grounding and pedagogical freedom}\label{sec:grounding-policy}
Table~\ref{tab:grounding-policy} separates source-required knowledge from permitted instructional invention. This policy controls generation and defines the denominator of the source-support assessment.
\begin{table}[H]\centering\small\caption{Grounding policy for generated teaching material.}\label{tab:grounding-policy}
\begin{tabularx}{\linewidth}{|p{3.0cm}|Y|}\hline
Part & Grounding rule\\\hline
Lesson content & Facts, definitions and formulas come from supplied passages with \code{[S\#]} citations. The model may paraphrase and connect ideas, but adds no new knowledge.\\\hline
Illustrative examples & Worked examples and code may be created, labelled ``Illustrative example,'' and consistent with source concepts. Code is checked by execution in an isolated environment where possible.\\\hline
Practice and Quiz & New scenarios may elicit Apply/Analyze processes. Correct answers must be derivable from cited sources. Distractors may express model-generated common misconceptions and need no citation; each explanation cites the passage supporting the correct answer.\\\hline
Source gaps & Missing knowledge required by an objective is flagged to the lecturer instead of supplied from model knowledge.\\\hline
Support assessment & Assess source-required facts, correct answers and explanations. Exclude explicitly illustrative material and distractors from the support denominator.\\\hline
\end{tabularx}\end{table}
Source support and pedagogical quality are evaluated separately, so a configuration that asks only easy factual questions cannot win on support alone\citep{aqgreview}.
\subsubsection{Slides, example, questions and explanation}
Slides and Example are generated from the Content as structured objects. Each Practice and Quiz question retrieves its own passages for one sub-objective. A question writes its reasoning before its choices, then the explanation drawn from that reasoning; the Explanation component collects them, so it costs no extra call. Choices are checked for a single correct index and rebalanced when the correct answer is visibly the longest. Learner-facing text uses the chosen output language, English or Vietnamese.
\subsubsection{Self-check and comparison condition}
Content and each question have one self-check by the model, with one regeneration on failure; the results are stored with the module. Schema validation and citation checks always run. The no-RAG control of the evaluation runs the same steps with the source block removed (Section~\ref{sec:conditions}); control modules cannot be published.
\subsubsection{Concurrency within a module}
Question generation took about 75\% of a module's time on a real model (Section~\ref{sec:results-latency}). Questions are therefore generated five at a time (\code{QUESTION_CONCURRENCY}), and each worker process runs four jobs at once. Rate-limited model calls are retried individually by background jobs, as described in Section~\ref{sec:architecture}; every stage records its duration, model calls and tokens, which produced the timing tables of Section~\ref{sec:results}.
\subsection{Gateway routing and model control}
The gateway holds the provider keys and maps aliases to models (Table~\ref{tab:model-aliases}). In normal use an alias may fall back to another model after errors; experiment mode disables fallback, and every call records the alias requested and the model that served it.
\begin{table}[H]\centering\small\caption{Gateway aliases and their use in this project.}\label{tab:model-aliases}
\begin{tabularx}{\linewidth}{|p{2.6cm}|p{3.8cm}|Y|}\hline
Alias & Model & Used for\\\hline
\code{claude-fast} & Claude Haiku 5.5 & Generation latency and load tests\\\hline
\code{openai-fast} & GPT-6 Luna & Load tests\\\hline
\code{claude-pro}, \code{openai-pro} & Claude Sonnet 5.5, GPT-6.1 Sol & Planned rubric judges\\\hline
\code{gemini-fast} & Gemini 3.5 Flash-Lite & Chat evaluation, module pilot, formula re-reading\\\hline
\code{groq-pro} & GPT-OSS-120B on Groq & Load control, Ragas and claim judge\\\hline
\code{groq-qwen} & Qwen 3.8 27B on Groq & Preliminary rubric judge\\\hline
\code{transcribe} & Whisper large-v3 turbo on Groq & Transcription\\\hline
\code{embed-default} & Qwen3-Embedding-0.6B, local & Embedding\\\hline
\end{tabularx}\end{table}
